\documentclass[12pt]{article}

\usepackage[a4paper, total={6in, 8in}]{geometry}
\usepackage{textcomp}

\begin{document}
\setlength{\parindent}{0pt}

\def \t {\textrightarrow}
\def \v {\vspace{1em}}
\def \bi {\begin{itemize}}
\def \ei {\end{itemize}}
\def \s[#1] {\section*{#1}}
\def \ss[#1] {\subsection*{#1}}
\def \sss[#1] {\subsubsection*{#1}}

\s[Costituzione]
La costituzione si compone di 12 principi fondamentali, che sono il riassunto di tutto il testo \t poi diritti e doveri, poi ordinamento della repubblica, poi varie \\
La parte sull ordinamento è divisa in titoli \\
\bi
\item rapporti civili
\item rapporti etico-sociali
\item rapporti 
\ei
Mentre ordinamento si parla di parlamento, etc. \\
La divisione della costituzione corrisponde al lavoro della assemblea costituente \t che aveva 156 deputati \t i deputati hanno nominato commissione di 75 deputati, chiamata commissione per costituzione, suddivisa in 3 sottocommissioni \\
Ciascuna si dedica a un argomento:
\bi
\item la prima: diritti e doveri dei cittadini
\item la seconda: organizzazione istituzionale dello stato
\item la terza: lineamenti economici e sociali
\ei
Ognuna ha redatto testo, poi riuniti in una bozza da 18 deputati \t da commissione nominata nel novembre del 46 \t poi bozza viene discussa da tutta l assemblea \\
Testo finale non viene votato da tutti \t ha avuto 62 voti contrari, che erano i voti dei monarchici, della destra conservatrice e del movimento qualunquista (il partito dell uomo qualunque) \\
Testo definitivo viene poi inviato a Enrico De Nicola \t la costituzione viene promulgata (ovver un atto formale con cui capo di stato valida una norma) \t il 27 dicembre del 47 viene promulgata la costituzione (validata quindi) \\
1 gennaio del 48 entra in vigore

\v

Le parti che la compongono sono orignate da come è stata creata, dalle 3 subcommissioni \\
"L'italia è una repubblica democratica fonadata sul lavoro, e la sovranita appartiene al popolo" \t la regolazione del lavoro è un dovere dello stato \\
Mentre l italia è una repubblica costituzionale \t è una democrazia, e potere del popolo deve stare dentro la costituzione \\
Sovranita si esercita in:
\bi
\item suffragio universale
\item referendum, forma di democrazia diretta (la nostra è rappresentativa) \t mentre con referendum cittadini sono chiamati ad esprimersi direttamente
\ei
Ci sono referendum 
\bi
  \item abrogativi
  \item consultivi
  \item confermativo di approvazione \t relativo alla modifica di una norma costituzionale \t es. l ultimo sulla magistratura
\ei
Non tutte le leggi pero possono essere soggete a referendum abrogativo, come:
\bi
  \item le leggi di bilancio
  \item quelle fiscali
  \item la ratifica di trattati internazionali
  \item quelle di amnistia (il parlamento concede un atto di clemenza nei confronti di qualcuno che ha commesson un reato \t viene cancellato pena e reato) e indulto (qua viene invece condonata soltanto la pena, il reato rimane nella fedina)
\ei

\v

Secondo articolo \t "La Repubblica riconosce e garantisce i diritti inviolabili dell'uomo, sia come singolo sia nelle formazioni sociali ove si svolge la sua personalità, e richiede l'adempimento dei doveri inderogabili di solidarietà politica, economica e sociale." \\
Definito articolo dei diritti e dei doveri \t perche vengono classificate due liberta:
\bi
  \item positive \t liberda di associazione, stampa, parola etc.
  \item liberta da arresti, etc.
\ei
Sono inviolabili, appartengono alla natura dell uomo \t non possono essere modificati, neanche con una revisione costituzionale \\
Bisogna pero salvaguardare l ordine pubblico \t quindi sono posti dei limiti, dei doveri \t il diritto va originarialmente limitato, per garantire ordine pubblico

\v

Parte prima - Diritti e doveri dei cittadini:
\bi
  \item titolo I - Rapporti civili
  \item titolo II - Rapporti etico-sociali
  \item titolo III - Rapporti economici
  \item titolo IV - Rapporti politici
\ei

\v

Terzo articolo \t "Tutti i cittadini hanno pari dignità sociale e sono eguali davanti alla legge, senza distinzione di sesso, di razza, di lingua, di religione, di opinioni politiche, di condizioni personali e sociali. È compito della Repubblica rimuovere gli ostacoli di ordine economico e sociale, che, limitando di fatto la libertà e l'eguaglianza dei cittadini" \\
Uguaglianza formale = trattare tutti allo stesso modo \t ma il problema è che io devo denifire cosa vuol dire trattare le persone in modo uguale, e come lo stato deve fare \\
Qua viene riconosciuta l uguaglianza, ma lo stato deve intervenire perche l uguaglianza sia reale \\
All inizio viene sancita un uguaglianza ideale, mentre nella seconda si dice che l uguaglianza deve essere nel concreto della vita del singolo ed e compito della repubblica \\
Questo non vuol dire che la costituzione non riconosce le differneze \t riconosciuta quella giuridica, ma poi c e riconoscimento necessario di differenze \t es. il riconoscimento delle minoranze linguistiche \\
Riconoscendo differenze, garantisco l'uguaglianza \t se si ignorano differenze (come minoranze), non vengono tutelate

\v

Articolo 4 \t "La Repubblica riconosce a tutti i cittadini il diritto al lavoro" \t diritto al lavoro, che è un diritto sociale \t il lavoro definisce il ruolo sociale (non è il censo, non una classe) \t il ruolo sociale nasce dal lavoro \\
Viene sancita identica protezione a tutti i tipi di lavoro \t il lavoro in quanto tale è protetto dalla costituzione \\
Articoli 35 e 37 sono dedicati al lavoro, mentre nel 38 si parla delle protezioni del lavoro attuate dallo stato \\
Diritto ai sindacati, e anche associazioni degli imprenditori (es confindustria etc.) \\
"Ogni cittadino ha il dovere di svolgere, secondo le proprie possibilità e la propria scelta, un'attività" \\
Articolo 40 sancisce il diritto di sciopero \t è ammesso lo sciopero di natura economica

\v

Articolo 5 \t "La Repubblica, una e indivisibile, riconosce e promuove le autonomie locali" \t ma repubblica è unica e indivisibile, non frammentaria con poteri locali \t pero riconosciute minoranze \\
Questo è antifascista \t l unita nazionale non è una costrizione, non è una distruzione delle minoranze etc (cosa che aveva fatto il fascismo) \\
Qua dichiarato invece un principio regionalistico \t sono riconusciute delle autonomie nel territorio

\v

Articolo 7 e 8 sono legati \\
Articolo 8 ribadisce l indipendenza tra stato e chiesa \t garantendo liberta a tutte le professioni, allora lo stato non ne tutela nessuna \t non è affare dello stato, che è completamente laico \\
I patti lateranensi prevedavano che lo stato si facesse carico della chiesa anche economicamente \t era prevista una congrua, ovvero lo stato provvedeva a una parte del reddito del clero \t con l 84 viene abolito, ma viene permesso l otto per mille (parte delle tasse puo essere contribuito alla chiesa, o in realta a qualunque religione) \\
Esistono in realta concordati anche con altre professioni \t es. con i protestanti nelle varie declinazioni (es. matrimoni protestanti hanno stesso valore di quelli cattolici) \t anche con la comunita ebraica

\v

Articolo 9 \t "La Repubblica promuove lo sviluppo della cultura e la ricerca scientifica e tecnica." \t è una premessa agli articoli 33 e 34 sulla scuola etc. \\
Questo è un articolo orientativo, non operativo \\
A livello di istruzione, non ci sono delle dichiarazioni operative (come per il lavoro) \t è piu orientativo \\
Vengono pero sancite:
\bi
  \item la liberta dell insegnamento \t ogni insegnante è libero di insegnare come crede, e che lo stato permette di istuire scuole non statali (pur che queste scuole private non enorino sul bilancio dello stato) \t questo pero ha diverse interpretazioni:
    \bi
      \item quindi stato non deve finanziare
      \item pero si possono supportare le famiglie che vogliono mandare i figli in scuola privata, stato pagava per esempio una parte della quota
    \ei
  \item le scuole paritarie devono attenersi a qud
\ei
Articolo 34 è l articolo del diritto allo studio \t tutti devono porter avere le stesse opportunita per studiare \\
"Tutela il paesaggio e il patrimonio storico e artistico della Nazione." \\
Nel 74 viene creato il ministero per i beni e attivita culturali ? \t prima era assorbito in parte dall istruizione, in parte dal miniestero degli interni \\
Ne 1902 e nel 1909 sono state fatte delle leggi che doevano tutelare i beni sul suolo italiano \t la legge Bottai (ministro educazione del 39) che mette le basi per quella che sara la vera tutela dei beni cultuarli in italia, che puoi confluisce nel codice dei beni culturali del 2004

\v

Articolo 10 \t "L'ordinamento giuridico italiano si conforma alle norme del diritto internazionale generalmente riconosciute." + "lo straniero ha diritto di asilo" \\
Anche qua di vede radisce anitfascista \t contro le posizioni xenofobe e imperialistiche dell fascismo, dove lo straniero veniva anzi perseguitato \\

\v

Articolo 11 \t "L'Italia ripudia la guerra come strumento di offesa" \\
"Limitazioni di sovranità necessarie ad un ordinamento che assicuri la pace e la giustizia fra le Nazioni" \t si vuole garantire la pace internazionale \t dove viene messa in discussione la pace a livello internazionale è consentito (non necessario) che ci siano delle limitazioni al potere sovrano, se si deve garantire pace \\
Qua si sente l eco della guerra \t si vuola aprire un epoca di cooperazione, con la costruzione di organismi sovranazionale \t questo articolo e il 10 sono il nucelo internazionalista della costituzione \\
Nell articolo 10 compare il diritto internazionale, secondo il quale la legislazione in merito alla tutela dello straniero si deve uniformare a quello internazionale \\
Viene sancito il diritto d asilo, il divieto di estradizione per motivi politici \\
Cittadini extracomunitari non appartengono all unione europea, e in italia sono soggetti a restrizioni maggiori \\
Le condizione per il diritto d asilo quali sono? \t è chi è al ???????????? \\
Italia entra nell onu nell 55

\v

Aritcolo 12 sulla bandiera \t il tricolore nasce nel 1797, ed è la bandiera della repubblica cis padana \t perche i giacobini si ispriano alla bandiera francese \\
Delle repubbliche giacobine la maggioranza non resiste, mentre quando napoelone crea regno italico repubblica conflusicono \t regno rimarra tale fino a congresso di vienna \t bandiera diventa il simbolo della resistenza alla restaurazione \\
Sara la bandiera dei mazziniani, e nel 48 diventa emblema della riscossa nazionale \t poi diventa bandiera del regno di italia \\
Non c era pero definizione, erano 3 colori per cambiati \t solo nel 1925 viene definito, con lo stemma sabaudo e la corona \\
Oggi solo i 3 colori
\end{document}
